\chapter{Conjunts finits i cardinalitat finita}\label{cap:finits}

\begin{objectius}
\apartat{Prerequisits} Del capítol~\ref{cap:naturals}: l'objecte \(\N\), la inducció (amb i sense paràmetres), les equacions de Peano, l'aritmètica i l'ordre, amb la convenció que les propietats de l'ordre es comproven sobre punts. Del capítol~\ref{cap:valors}: sumes disjuntes, complements i distributivitat. Dels capítols~\ref{cap:imatges} i~\ref{cap:exponencials}: factorització imatge, exponencials i composició interna. L'axioma d'elecció només apareix en un punt, assenyalat.
\apartat{Objectius} En acabar el capítol, el lector hauria de poder:
\begin{itemize}
  \item distingir els ordinals finits \([k]\), definits al metanivell per a numerals estàndard, dels segments inicials \(\seg{n}\), que formen una família interna, i explicar per què la finitud es defineix amb els segons;
  \item construir el conjunt de cardinals \(\operatorname{Card}(A)\hookrightarrow\N\) codificant les bijeccions \(\seg{n}\to A\) dins un exponencial fix, i definir finitud i cardinal a partir d'ell;
  \item demostrar, per inducció interna sobre un predicat codificat, que un monomorfisme \(\seg{m}\hookrightarrow\seg{n}\) força \(m\le n\), i deduir-ne la unicitat del cardinal;
  \item demostrar el principi dels calaixos i l'equivalència entre mono, epi i iso en cardinal igual, i saber on intervé l'elecció;
  \item demostrar els recomptes \(|A+B|=|A|+|B|\), \(|A\times B|=|A|\cdot|B|\), \(|P(A)|=2^{|A|}\) i \(|B^A|=|B|^{|A|}\) per inducció.
\end{itemize}
\end{objectius}

Al llarg del llibre, la paraula \emph{finit} ha aparegut només com a
comentari extern informal: dèiem que una construcció produïa «un nombre
finit» de casos o que un diagrama tenia «finites» fletxes, sense que
aquesta finitud fos res precís. Ara que disposem de l'objecte dels nombres
naturals (capítol~\ref{cap:naturals}), podem convertir-la en una noció
\emph{precisa} i \emph{interna}: un conjunt \(A\) serà finit quan existeixi
un punt \(n\colon\mathbf1\to\N\) amb \(\seg{n}\approx A\), és a dir, quan
\(A\) sigui equipotent a algun dels segments inicials
(definició~\ref{def:segment-inicial}) que el successor genera. Aquesta és
una propietat interna, expressada dins el llenguatge mitjançant la família
uniforme dels segments \(n\mapsto\seg{n}\), i no requereix cap axioma nou.
Els ordinals finits \([k]\) que introduïm tot seguit
(definició~\ref{def:ordinal-finit}) no en són la definició, sinó només la
\emph{lectura externa} per als numerals estàndard
(observació~\ref{obs:ordinals-metanivell}): una notació de càlcul còmoda,
lligada als segments per la coincidència de
l'observació~\ref{obs:segment-numeral}. El capítol és, més aviat, una
primera aplicació extensa de tot el que s'ha acumulat ---sumes,
exponencials, conjunt de parts, factorització imatge i inducció--- i el
pont cap als exercicis computacionals de la part pràctica.

\section{Els ordinals finits}\label{sec:ordinals-finits}

La successió \(\mathbf1,\;\mathbf1+\mathbf1,\;\mathbf1+\mathbf1+\mathbf1,
\dots\) no és, de moment, res més que una família de conjunts descrita amb
punts suspensius. Comencem donant-li un nom precís \emph{al metanivell},
una recursió externa sobre els numerals estàndard; tot seguit
(secció~\ref{sec:segments-interns}) en construirem la versió \emph{interna}
i uniforme, els segments inicials, que és la que sostindrà tota la teoria.

\begin{definicio}[Ordinal finit]\index{Ordinal finit}\label{def:ordinal-finit}
Per a cada natural estàndard \(k\) del metallenguatge definim
\emph{externament} un conjunt \([k]\), l'\emph{ordinal finit}
associat a \(k\), per la recursió metateòrica
\[
[0]:=\mathbf0,
\qquad
[k+1]:=[k]+\mathbf1,
\]
on \(\mathbf0\) és el conjunt inicial (capítol~\ref{cap:aplicacions}),
\(+\) la suma de conjunts (definició~\ref{def:suma-conjunts}) i
\(\mathbf1\) el conjunt terminal. Aquesta notació assigna un conjunt
\([k]\), concret llevat d'isomorfisme, a cada numeral estàndard, un per a
cada \(k\) del metallenguatge; \emph{no} s'estén a un punt intern arbitrari
\(n\colon\mathbf1\to\N\). Per al pas de \([k]\) a \([k+1]\) escriurem
\(\iota_k\colon[k]\to[k+1]\) per a la injecció canònica \(i_0\) de la suma
\([k]+\mathbf1\), i \(\ast_k\colon\mathbf1\to[k+1]\) per a la injecció
\(i_1\) del sumand nou.
\end{definicio}

\begin{observacio}[Sobre l'estatut de la recursió d'ordinals]\label{obs:ordinals-metanivell}
La regla anterior és una recursió al metanivell, no una aplicació interna
cap a un «conjunt de tots els conjunts», que no existeix. A cada numeral
estàndard \(\bar k\colon\mathbf1\to\N\) ---el zero, el successor del zero,
i així successivament--- li correspon el conjunt \([k]\), concret llevat
d'isomorfisme, un per a cada \(k\) del metallenguatge. No pretenem que tot
punt global d'un \(\N\) intern sigui d'aquesta forma: en una teoria de
primer ordre no es pot identificar sense més una variable interna
\(n\colon\mathbf1\to\N\) amb un numeral construït externament, i per tant la
notació \([\,\cdot\,]\) no s'estén a un \(n\) intern arbitrari.

Per aquesta raó, la finitud i la cardinalitat internes \emph{no} es
basaran en la família metanivell \(k\mapsto[k]\), sinó en la família
\emph{interna} dels segments inicials \(n\mapsto\seg{n}\)
(definició~\ref{def:segment-inicial}), que prové d'un sol objecte, la
relació \(\operatorname{Lt}\hookrightarrow\N\times\N\). Els ordinals
\([k]\) es conservaran només com a \emph{notació de càlcul} per als
numerals estàndard, lligada als segments per la coincidència externa de
l'observació~\ref{obs:segment-numeral}.
\end{observacio}

Convé fixar de seguida les mides petites, que reapareixeran a cada
inducció.

\begin{proposicio}[Els primers ordinals]\label{prop:ordinals-petits}
Es tenen isomorfismes canònics \([1]\cong\mathbf1\) i
\([2]\cong\mathbf1+\mathbf1\cong\mathbf2\). A més, per a cada numeral
estàndard \(k\) la injecció \(\iota_k\colon[k]\hookrightarrow[k+1]\) és un
monomorfisme, i \(\ast_k\) és el seu complement dins \([k+1]\).
\end{proposicio}

\begin{prova}
El cas \(k=0\) de la definició dona \([1]=\mathbf0+\mathbf1\cong\mathbf1\)
per la neutralitat del conjunt inicial respecte de la suma
(capítol~\ref{cap:valors}); aplicant-ho un cop més,
\([2]=[1]+\mathbf1\cong\mathbf1+\mathbf1\), que és el conjunt de valors de
veritat \(\mathbf2\). Que \(\iota_k\) és un monomorfisme amb complement
\(\ast_k\) és la disjunció de la suma (definició~\ref{def:suma-disjunta}):
en una suma disjunta les dues injeccions són monomorfismes que
representen parts complementàries. Com que aquestes afirmacions són un
esquema sobre numerals estàndard, cada instància es comprova externament.
\end{prova}

\section{Els segments inicials com a família interna}\label{sec:segments-interns}

Els ordinals \([k]\) de la definició~\ref{def:ordinal-finit} pateixen la
limitació que l'observació~\ref{obs:ordinals-metanivell} ja assenyalava:
la notació \(k\mapsto[k]\) només està definida per als numerals estàndard
del metallenguatge i no és, per si sola, una família interna. Ara que
disposem de l'ordre de
\(\N\) com a \emph{part} \(\operatorname{Lt}\hookrightarrow\N\times\N\)
(definició~\ref{def:ordre-naturals}), podem donar una segona presentació
dels ordinals finits que \emph{sí} és interna i uniforme, i comprovar que
coincideix amb la primera. Aquesta presentació és la que farà de la
finitud una propietat interna.

\begin{definicio}[Segment inicial]\index{Segment inicial}\label{def:segment-inicial}
Per a cada punt \(n\colon\mathbf1\to\N\), el \emph{segment inicial}
\(\seg{n}\hookrightarrow\N\) és la fibra de la part
\(\operatorname{Lt}\hookrightarrow\N\times\N\) sobre \(n\): la part de
\(\N\) formada pels \(k\) amb \(k<n\), és a dir la imatge inversa
\[
\seg{n}:=\{\,k\colon\mathbf1\to\N \mid k<n\,\}
=\operatorname{Lt}\bigl|_{\N\times\{n\}} ,
\]
obtinguda com el producte fibrat de \(\operatorname{Lt}\hookrightarrow
\N\times\N\) al llarg de \(\langle\id_\N,n\circ{!}\rangle\colon\N\to
\N\times\N\) (definició~\ref{def:pullback}).
\end{definicio}

\begin{observacio}[La família de segments és interna i uniforme]\label{obs:segments-interns}
A diferència de \(k\mapsto[k]\), la regla \(n\mapsto\seg{n}\) \emph{sí} que és una única part interna: prové d'un sol
objecte, la relació \(\operatorname{Lt}\hookrightarrow\N\times\N\), de la
qual \(\seg{n}\) és la fibra sobre \(n\). Dit d'una altra
manera, \(\operatorname{Lt}\), vista com una aplicació
\(\operatorname{Lt}\to\N\) per la segona projecció, és la
\emph{família uniforme} \((\seg{n})_{n\in\N}\) de tots els
segments alhora, indexada internament per \(\N\)
(capítol~\ref{cap:relacions}). És aquesta uniformitat la que permetrà
quantificar sobre segments dins el llenguatge, cosa que amb \([k]\) sol no
era possible.
\end{observacio}

\begin{observacio}[Coincidència externa amb els numerals estàndard]\label{obs:segment-numeral}
Per a cada numeral estàndard \(k\) del metallenguatge, hi ha un isomorfisme
\(\seg{\bar k}\cong[k]\) entre el segment del numeral \(\bar
k\colon\mathbf1\to\N\) i l'ordinal finit \([k]\). Aquesta és una afirmació
\emph{externa}, un esquema amb una instància per a cada \(k\), i es
comprova per inducció metateòrica sobre \(k\):
\[
\seg{\bar 0}\cong\mathbf0=[0],
\qquad
\seg{\overline{k+1}}\cong\seg{\bar k}+\mathbf1\cong[k]+\mathbf1=[k+1],
\]
on la primera línia val perquè \(\{j:j<0\}\) és inicial, i el pas usa que
\(j<\overline{k+1}\) equival a \(j\le\bar k\), és a dir \(j<\bar k\) o
\(j=\bar k\) (proposició~\ref{prop:ordre-total}), cosa que descompon
\(\seg{\overline{k+1}}\) en la suma disjunta \(\seg{\bar k}+\mathbf1\).

No enunciem cap identificació \(\seg{n}\cong[n]\) per a un \(n\colon
\mathbf1\to\N\) \emph{intern} arbitrari: això requeriria una inducció
interna sobre un objecte, \([n]\), que no és intern. Els \([k]\) queden,
doncs, com a abreviatura còmoda per als numerals concrets, i tota
construcció interna del capítol es fa amb els segments \(\seg{n}\).
\end{observacio}

El que \emph{sí} és intern, i és tot el que la teoria de la finitud
necessita, és la manera com un segment creix en passar al successor. Ho
enunciem sobre una variable interna \(n\).

\begin{proposicio}[Descomposició successora del segment]\label{prop:segment-successor}
Per a tota variable interna \(n\colon\mathbf1\to\N\) hi ha un isomorfisme
\[
\seg{s\circ n}\ \cong\ \seg{n}+\mathbf1,
\]
en què el sumand \(\seg{n}\) és la part de \(\seg{s\circ n}\) dels
\(k<n\) i el sumand \(\mathbf1\) és el punt \(n\) mateix. A més
\(\seg{0}\cong\mathbf0\).
\end{proposicio}

\begin{prova}
\(\seg{0}=\{k:k<0\}\) és inicial, perquè \(k<0\) no es compleix per a cap
punt \(k\) (de \(k+s\circ j=0\) sortiria \(s\circ(k+j)=0\), contra el
fet que el zero no és cap successor, teorema~\ref{teo:peano}(2)).

Per al successor, un punt \(k\in\seg{s\circ n}\) compleix \(k<s\circ n\),
que equival a \(k\le n\), és a dir \(k<n\) o bé \(k=n\), i no totes dues
alhora perquè \(n\not<n\) (pas del successor de l'ordre,
proposició~\ref{prop:ordre-total}); com a tota la secció de l'ordre, el
raonament és sobre punts, i la separació global el converteix en una
igualtat de parts (observació~\ref{obs:ordre-sobre-punts}). Aquesta disjunció exclusiva
descompon la part \(\seg{s\circ n}\hookrightarrow\N\) en dues parts
complementàries: la dels \(k<n\), que és \(\seg{n}\), i la del punt únic
\(k=n\), que és una còpia de \(\mathbf1\). Una part i el seu complement
formen una suma disjunta (teorema~\ref{teo:part-complementari}), de manera
que \(\seg{s\circ n}\cong\seg{n}+\mathbf1\).
\end{prova}

\section{Equipotència i el conjunt de cardinals}\label{sec:equipotencia}

Per definir la finitud com una propietat \emph{interna} ---i, alhora, el
cardinal com una funció ben definida--- ens cal, per a un conjunt \(A\)
fixat, poder reunir en una part de \(\N\) tots els \(n\) tals que \(\seg{n}\)
és isomorf a \(A\). Aquesta part es construeix sense cap circularitat, ni
amb cap axioma nou, a partir de dues eines ja disponibles: el subconjunt
dels isomorfismes dins un exponencial, i la factorització imatge. La
construcció és el cas finit d'una noció més general ---l'\emph{equipotència}
de dos conjunts--- que reprendrem quan comparem cardinals infinits.

\begin{definicio}[El subconjunt dels isomorfismes]\index{El subconjunt dels isomorfismes}\label{def:iso-exponencial}
Siguin \(A\) i \(B\) dos conjunts. El \emph{subconjunt dels isomorfismes}
\(\operatorname{Isos}(A,B)\hookrightarrow B^{A}\) és la part de
l'exponencial \(B^{A}\) formada pels punts \(f\) que admeten una inversa
bilateral. Es construeix internament així. Per la composició interna
(proposició~\ref{prop:composicio-interna}) hi ha aplicacions
\[
c_{A}\colon A^{B}\times B^{A}\to A^{A},
\qquad
c_{B}\colon B^{A}\times A^{B}\to B^{B},
\]
que sobre punts envien \((g,f)\mapsto g\circ f\) i \((f,g)\mapsto f\circ
g\) respectivament, i identitats internes \(\jmath_A\colon\mathbf1\to
A^{A}\), \(\jmath_B\colon\mathbf1\to B^{B}\). Considerem l'aplicació
\[
\Phi\colon B^{A}\times A^{B}\longrightarrow A^{A}\times B^{B},
\qquad
(f,g)\longmapsto\bigl(c_A(g,f),\,c_B(f,g)\bigr)=(g\circ f,\ f\circ g),
\]
i l'equalitzador (definició~\ref{def:equalitzador}) de \(\Phi\) amb
l'aplicació constant de valor \((\jmath_A,\jmath_B)\): la part
\(P\hookrightarrow B^{A}\times A^{B}\) dels parells \((f,g)\) amb \(g\circ
f=\id_A\) i \(f\circ g=\id_B\). Aleshores \(\operatorname{Isos}(A,B)\) és la
imatge (teorema~\ref{teo:factoritzacio}) de \(P\) per la projecció a
\(B^{A}\): els \(f\) per als quals existeix un tal \(g\).
\end{definicio}

\begin{observacio}[Els punts de \(\operatorname{Isos}(A,B)\)]\label{obs:iso-punts}
Pels punts globals de l'exponencial (teorema~\ref{teo:punts-exponencial}),
un punt \(\mathbf1\to\operatorname{Isos}(A,B)\) és exactament una aplicació
\(f\colon A\to B\) que és un isomorfisme (definició~\ref{def:isomorfisme}).
Així, \(\operatorname{Isos}(A,B)\) és no inicial si i només si existeix un
isomorfisme \(A\cong B\). La part es defineix per a tot parell de conjunts
amb un únic bloc d'operacions internes, sense triar cap isomorfisme.
\end{observacio}

\begin{definicio}[Equipotència]\index{Equipotència}\label{def:equipotencia}
Dos conjunts \(A\) i \(B\) són \emph{equipotents}, i escrivim \(A\approx
B\), si \(\operatorname{Isos}(A,B)\) és no inicial ---equivalentment, si
existeix un isomorfisme \(A\cong B\). La relació \(\approx\) és reflexiva
(la identitat), simètrica (la inversa d'un isomorfisme) i transitiva (la
composició), és a dir, una relació d'equivalència. Cal entendre aquesta
afirmació \emph{metateòricament}: com que no hi ha cap conjunt de tots els
conjunts, \(\approx\) no és una part interna d'un producte, sinó una
relació entre conjunts formulada al metanivell; reflexivitat, simetria i
transitivitat s'hi comproven exhibint, en cada cas, l'isomorfisme
corresponent.
\end{definicio}

Ara fixem \(A\) i fem variar el segment. La família uniforme
\((\seg{n})_{n\in\N}\) (observació~\ref{obs:segments-interns}) permet
reunir en un sol objecte tots els isomorfismes \(\seg{n}\to A\), i
projectar-los sobre \(\N\) dona la part buscada.

\begin{definicio}[El conjunt de cardinals d'un conjunt]\index{El conjunt de cardinals d'un conjunt}\label{def:card-A}
Sigui \(A\) un conjunt. Cada segment \(\seg{n}\) és una part de \(\N\)
(definició~\ref{def:segment-inicial}), de manera que una aplicació
\(\varphi\colon\seg{n}\to A\) s'estén de manera única a una aplicació
\(\bar\varphi\colon\N\to A+\mathbf1\) que val \(\varphi\) sobre \(\seg{n}\)
i el punt afegit \(\ast\colon\mathbf1\to A+\mathbf1\) fora de \(\seg{n}\)
(definició per casos amb el complement, capítol~\ref{cap:valors}). Així
totes les aplicacions des de tots els segments viuen dins un \emph{únic}
exponencial fix, \((A+\mathbf1)^{\N}\), i podem reunir-les en un sol
objecte.

Definim una part d'un objecte \emph{fix}, sense exponencials de domini
variable:
\[
\mathcal E_A\ \hookrightarrow\ \N\times(A+\mathbf1)^{\N}\times\N^{A},
\]
els punts de la qual són els triples \((n,g,h)\) ---amb
\(g\colon\N\to A+\mathbf1\) i \(h\colon A\to\N\)--- que satisfan
internament les tres condicions següents, on \(i_A\colon A\to A+\mathbf1\)
i \(i_1\colon\mathbf1\to A+\mathbf1\) són les injeccions de la suma:
\begin{enumerate}
\item[(i)] per a tot \(k<n\), existeix (un únic) \(a\) amb \(g(k)=i_A(a)\)
i \(h(a)=k\);
\item[(ii)] per a tot \(a\in A\), \(h(a)<n\) i \(g(h(a))=i_A(a)\);
\item[(iii)] per a tot \(k\) amb \(k\not<n\), \(g(k)=i_1(\ast)\).
\end{enumerate}
Aquestes condicions només fan servir avaluacions dels exponencials
\emph{fixos} \((A+\mathbf1)^{\N}\) i \(\N^{A}\), l'ordre \(<\)
(definició~\ref{def:ordre-naturals}), igualtats i quantificadors acotats,
tots ja construïts; \(\mathcal E_A\) n'és la part corresponent, tallada per
equalitzadors i imatges inverses. Les condicions (i)--(ii) diuen que
\(g\) i \(h\) exhibeixen una bijecció entre \(\seg{n}\) i \(A\) (cadascuna
inversa de l'altra sobre el seu domini), i (iii) fixa \(g\) fora de
\(\seg{n}\), de manera que un tal triple existeix si i només si
\(\seg{n}\approx A\).

El \emph{conjunt de cardinals} de \(A\) és la imatge
(teorema~\ref{teo:factoritzacio}) de la projecció
\(\pi\colon\mathcal E_A\to\N\):
\[
\operatorname{Card}(A)\ :=\ \operatorname{Im}(\pi)\ =\
\{\,n\colon\mathbf1\to\N \mid \seg{n}\approx A\,\}
\ \hookrightarrow\ \N .
\]
\end{definicio}

\begin{proposicio}[\(\operatorname{Card}(A)\) és una part interna]\label{prop:card-interna}
Per a tot conjunt \(A\), \(\operatorname{Card}(A)\) és una part ben
definida de \(\N\), i un punt \(n\colon\mathbf1\to\N\) hi pertany si i
només si \(\seg{n}\approx A\). La construcció no fa servir cap axioma
d'elecció ni cap noció prèvia de finitud o cardinal.
\end{proposicio}

\begin{prova}
La part \(\mathcal E_A\hookrightarrow\N\times(A+\mathbf1)^{\N}\times\N^{A}\)
és interna per construcció (definició~\ref{def:card-A}): les condicions
(i)--(iii) són conjuncions d'equacions entre avaluacions d'exponencials
fixos i de condicions acotades per l'ordre \(<\), totes tallables amb
equalitzadors, imatges inverses i quantificadors acotats ja disponibles
(capítols~\ref{cap:exponencials} i~\ref{cap:parts-interna}). La imatge
d'una aplicació existeix sempre (teorema~\ref{teo:factoritzacio}), de
manera que \(\operatorname{Card}(A)=\operatorname{Im}(\pi)\) és una part de
\(\N\).

Un punt \(n\) es factoritza per aquesta imatge si i només si la fibra
\((\mathcal E_A)_n\) és no inicial (exhaustivitat de la coprojecció sobre
la imatge, capítol~\ref{cap:imatges}). Ara bé, un punt de \((\mathcal
E_A)_n\) és un parell \((g,h)\) que satisfà (i)--(iii); per (i) i (ii),
\(g\) restringida a \(\seg{n}\) i \(h\) són mútuament inverses, i
defineixen un isomorfisme \(\seg{n}\to A\), és a dir \(\seg{n}\approx A\).
Recíprocament, un isomorfisme \(\varphi\colon\seg{n}\to A\) dona el triple
\((n,\bar\varphi,\varphi^{-1})\), amb \(\bar\varphi\) l'extensió per
\(\ast\), que satisfà (i)--(iii). Així \(n\in\operatorname{Card}(A)\) si i
només si \(\seg{n}\approx A\). Cap pas no ha usat la finitud ni el
cardinal, sinó només exponencials fixos, l'ordre, equalitzadors i la
imatge.
\end{prova}

\begin{definicio}[Conjunt finit i cardinal]\index{Conjunt finit i cardinal}\label{def:conjunt-finit}
Un conjunt \(A\) és \emph{finit} si \(\operatorname{Card}(A)\) és no
inicial (definició~\ref{def:card-A}); és a dir, si existeix un punt
\(n\colon\mathbf1\to\N\) amb \(\seg{n}\approx A\). Quan \(A\) és finit,
provarem tot seguit que \(\operatorname{Card}(A)\) té un \emph{únic} punt
(teorema~\ref{teo:unicitat-cardinal}); aquest punt és el \emph{cardinal}
de \(A\), que denotarem \(|A|\colon\mathbf1\to\N\).
\end{definicio}

\begin{observacio}[La finitud és una propietat interna]\label{obs:finitud-propietat}
Aquesta definició \emph{sí} és interna, a diferència del que passaria amb
\([n]\). La clau és que \(\operatorname{Card}(A)\) s'ha construït
(definició~\ref{def:card-A}, proposició~\ref{prop:card-interna}) com una
part genuïna de \(\N\) ---la imatge d'una projecció---, sense recórrer a
cap noció prèvia de finitud ni de cardinal, sinó només a la família interna
de segments (observació~\ref{obs:segments-interns}), al subconjunt dels
isomorfismes d'un exponencial (definició~\ref{def:iso-exponencial}) i a la
factorització imatge (teorema~\ref{teo:factoritzacio}). Per tant «\(A\) és
finit» ---que \(\operatorname{Card}(A)\) sigui no inicial--- és una
propietat interna de \(A\), i no un esquema indexat per numerals com ho era
la caracterització amb \([n]\).

No hi ha circularitat entre finitud i cardinal: primer es construeix la
part \(\operatorname{Card}(A)\); després, «finit» n'afirma la no-inicialitat
i el «cardinal» n'és l'únic punt. Cap de les dues nocions no s'ha usat en
la construcció de \(\operatorname{Card}(A)\). Que aquest punt és únic ---i
per tant que \(|A|\) és una funció ben definida sobre els conjunts finits,
invariant per isomorfisme--- és el contingut del
teorema~\ref{teo:unicitat-cardinal}, que provem tot seguit.
\end{observacio}

\begin{observacio}[Relació amb la finitud estàndard]\label{obs:finitud-equivalencia}
La definició interna es pot comparar, \emph{externament}, amb la noció
ingènua de finitud «isomorf a algun \([k]\)». Per a cada numeral estàndard
\(k\), la coincidència \(\seg{\bar k}\cong[k]\)
(observació~\ref{obs:segment-numeral}) dona: si \(A\cong[k]\) per a algun
numeral \(k\), aleshores \(\seg{\bar k}\approx A\), de manera que
\(\bar k\in\operatorname{Card}(A)\) i \(A\) és finit en el sentit intern
(definició~\ref{def:conjunt-finit}). Recíprocament, si es disposa
\emph{externament} d'un punt global \(n\in\operatorname{Card}(A)\) que
sigui un numeral estàndard \(\bar k\), aleshores \(A\cong\seg{\bar
k}\cong[k]\).

Convé distingir tres nivells, que no s'han de confondre:
\begin{enumerate}
\item \(\operatorname{Card}(A)\hookrightarrow\N\) és una \emph{part
interna}, sempre ben definida (proposició~\ref{prop:card-interna});
\item «\(A\) és finit» és l'afirmació, \emph{externa}, que aquesta part és
no inicial ---és a dir, que té algun punt global;
\item un tal punt global és un \(n\colon\mathbf1\to\N\) que no té per què
ser un numeral estàndard; la identificació amb un \([k]\) concret només és
possible quan, externament, se sap que \(n\) és un numeral.
\end{enumerate}
Els resultats interns del capítol (comparació, unicitat, recomptes) no
requereixen cap identificació amb numerals: es demostren sobre els segments
\(\seg{n}\) amb \(n\) intern.
\end{observacio}

\begin{observacio}[Treballem sobre els segments interns]\label{obs:seg-vs-corxet}
En tot el que segueix, els resultats interns ---comparació de mides,
unicitat del cardinal, recomptes--- es demostren \emph{sobre els segments
\(\seg{n}\)} amb \(n\colon\mathbf1\to\N\) intern, i les inducions es fan
amb el principi d'inducció (teorema~\ref{teo:induccio}) usant la
descomposició successora \(\seg{s\circ n}\cong\seg{n}+\mathbf1\)
(proposició~\ref{prop:segment-successor}) allà on abans s'hauria escrit
\([s\circ n]=[n]+\mathbf1\). Així cap prova no depèn de la família
metanivell \([n]\): la finitud i el cardinal són interns
(definició~\ref{def:conjunt-finit}), i els \([k]\) queden reservats a la
lectura externa amb numerals concrets (observació~\ref{obs:segment-numeral}).
\end{observacio}

El nucli de tota la teoria és el fet que dos ordinals isomorfs són iguals.
El demostrarem per inducció (lema~\ref{lema:comparacio-ordinals}), i en derivarem alhora la unicitat del
cardinal i el principi dels calaixos.

Abans, però, cal assegurar que la inducció d'aquest lema \emph{sí} que és
interna. La dificultat és la mateixa que resolguérem per a
\(\operatorname{Card}(A)\): tant el domini \(\seg{m}\) com el codomini
\(\seg{n}\) varien amb els punts \(m,n\), de manera que «existeix un
monomorfisme \(\seg{m}\hookrightarrow\seg{n}\)» no és, escrit així, una
fórmula amb \(m\) lliure sobre la qual aplicar el principi d'inducció
(teorema~\ref{teo:induccio}). Ho resolem codificant tots els monomorfismes
entre segments dins un \emph{exponencial fix}, exactament com fèrem amb
\(\mathcal E_A\).

\begin{definicio}[Codificació uniforme dels monomorfismes entre segments]\index{Codificació uniforme dels monomorfismes entre segments}\label{def:mono-segments}
De la mateixa manera que a la definició~\ref{def:card-A}, una aplicació
\(\seg{m}\to\seg{n}\subseteq\N\) s'estén de manera única a una aplicació
\(\N\to\N+\mathbf1\) que hi val el punt afegit fora de \(\seg{m}\). Definim
una part d'un objecte \emph{fix},
\[
\mathcal M\ \hookrightarrow\ \N\times\N\times(\N+\mathbf1)^{\N},
\]
els punts \((m,n,g)\) de la qual ---amb \(g\colon\N\to\N+\mathbf1\), i
\(i_\N\colon\N\to\N+\mathbf1\), \(i_1\colon\mathbf1\to\N+\mathbf1\) les
injeccions de la suma--- satisfan internament:
\begin{enumerate}
\item[(m1)] per a tot \(k<m\), existeix un únic \(\ell<n\) amb
\(g(k)=i_\N(\ell)\);
\item[(m2)] per a tot \(k\) amb \(k\not<m\), \(g(k)=i_1(\ast)\);
\item[(m3)] per a tot \(k,k'<m\), si \(g(k)=g(k')\) aleshores \(k=k'\).
\end{enumerate}
Aquestes condicions només fan servir avaluacions de l'exponencial \emph{fix}
\((\N+\mathbf1)^{\N}\), l'ordre \(<\) (definició~\ref{def:ordre-naturals}),
igualtats i quantificadors acotats; \(\mathcal M\) n'és la part
corresponent, tallada per equalitzadors i imatges inverses. Les condicions
(m1)--(m2) diuen que \(g\) codifica una aplicació \(\seg{m}\to\seg{n}\)
(estesa per \(\ast\)), i (m3) que és injectiva. Així la fibra
\(\mathcal M_{m,n}\) sobre \((m,n)\) és no inicial si i només si existeix un
monomorfisme \(\seg{m}\hookrightarrow\seg{n}\).
\end{definicio}

Amb aquesta codificació, «existeix un monomorfisme
\(\seg{m}\hookrightarrow\seg{n}\)» esdevé la condició interna
«\(\mathcal M_{m,n}\) és no inicial», i podem formar el predicat intern
\[
P(m)\ :\Longleftrightarrow\ \forall n\colon\mathbf1\to\N\ \bigl(\mathcal
M_{m,n}\text{ és no inicial}\ \Rightarrow\ m\le n\bigr),
\]
que sí que és una fórmula amb \(m\) com a única variable lliure, apta per al
principi d'inducció (teorema~\ref{teo:induccio}). La prova del lema aplica
la inducció interna precisament a \(P\).

\begin{lema}[Comparació de segments]\label{lema:comparacio-ordinals}
Siguin \(m,n\colon\mathbf1\to\N\).
\begin{enumerate}
\item Si existeix un monomorfisme \(\seg{m}\hookrightarrow\seg{n}\),
aleshores \(m\le n\).
\item Si existeix un isomorfisme \(\seg{m}\cong\seg{n}\), aleshores \(m=n\).
\end{enumerate}
\end{lema}

\begin{prova}
La segona afirmació se segueix de la primera: un isomorfisme dona
monomorfismes \(\seg{m}\hookrightarrow\seg{n}\) i
\(\seg{n}\hookrightarrow\seg{m}\), d'on \(m\le n\) i \(n\le m\), i per
l'antisimetria (proposició~\ref{prop:ordre-total}) \(m=n\).

Provem la primera aplicant el principi d'inducció
(teorema~\ref{teo:induccio}) al predicat intern \(P(m)\) definit més amunt
(definició~\ref{def:mono-segments}); com que \(P\) quantifica internament
sobre \(n\), la inducció és uniforme en \(n\) sense sortir del llenguatge.
Per a \(P(0)\): si \(\mathcal M_{0,n}\) és no inicial, aleshores \(0\le n\)
sempre (proposició~\ref{prop:ordre-total}). Per al pas, suposem \(P(m)\) i
provem \(P(s\circ m)\): sigui \(n\) tal que \(\mathcal M_{s\circ m,\,n}\)
és no inicial, és a dir hi ha un monomorfisme
\(f\colon\seg{s\circ m}\hookrightarrow\seg{n}\). Com que \(\seg{s\circ m}\cong\seg{m}+\mathbf1\)
(proposició~\ref{prop:segment-successor}) és no inicial, també ho és
\(\seg{n}\); per tant \(n\ne0\) ---si \(n=0\), \(\seg{0}\cong\mathbf0\) no rebria cap
aplicació des d'un conjunt no inicial (estrictesa de \(\mathbf0\),
capítol~\ref{cap:valors})---, és a dir \(n=s\circ n'\) per a un (únic)
punt \(n'\) (lema~\ref{lema:zero-o-successor}), i
\(\seg{n}\cong\seg{n'}+\mathbf1\)
(proposició~\ref{prop:segment-successor}).

Escrivim \(\seg{s\circ m}\cong\seg{m}+\mathbf1\) amb les injeccions
\(\iota\colon\seg{m}\hookrightarrow\seg{s\circ m}\) i
\(\ast\colon\mathbf1\to\seg{s\circ m}\), i
\(\seg{n}\cong\seg{n'}+\mathbf1\) amb
\(\jmath\colon\seg{n'}\hookrightarrow\seg{n}\) i
\(\ast'\colon\mathbf1\to\seg{n}\). Volem un monomorfisme
\(g\colon\seg{m}\hookrightarrow\seg{n'}\). El pas essencial és decidir on
va a parar el punt nou \(f\circ\ast\colon\mathbf1\to\seg{n}\). Com que
\(\seg{n}\cong\seg{n'}+\mathbf1\) és una suma disjunta, aquest punt global
es factoritza per exactament una de les dues injeccions
(teorema~\ref{teo:exhaustivitat-sumes}: tot punt global d'una suma prové
d'un únic sumand); distingim els dos casos.

\emph{Cas 1: \(f\circ\ast=\ast'\).} Aleshores la imatge inversa
\((f\circ\iota)^{-1}(\ast')\hookrightarrow\seg{m}\) és buida: un punt
\(x\) seu compliria \(f\circ\iota\circ x=\ast'=f\circ\ast\), i la
injectivitat de \(f\) donaria \(\iota\circ x=\ast\), impossible perquè
\(\iota\) i \(\ast\) són parts complementàries de \(\seg{s\circ m}\) (la
seva intersecció és inicial i no té punts); sense punts, la imatge
inversa és inicial pel lema del buit (lema~\ref{lema:buit}). Per tant
\(f\circ\iota\) es factoritza pel complement de \(\ast'\) (descomposició
\(\seg{n}\cong\ast'+(\ast')^c\) i estabilitat), que és
\(\jmath\colon\seg{n'}\hookrightarrow\seg{n}\):
existeix un únic \(g\colon\seg{m}\to\seg{n'}\) amb \(\jmath\circ
g=f\circ\iota\). Aquest \(g\) és un monomorfisme perquè \(\jmath\circ g\)
ho és i \(\jmath\) és un monomorfisme.

\emph{Cas 2: \(f\circ\ast=\jmath\circ q\)} per a un únic punt
\(q\colon\mathbf1\to\seg{n'}\). Reduïm al Cas~1 permutant a \(\seg{n}\) el
punt \(\jmath\circ q\) amb el punt nou \(\ast'\). Definim
\(\tau\colon\seg{n}\to\seg{n}\) \emph{per casos}, usant la descomposició de
\(\seg{n'}\) en el singleton determinat per \(q\) i el seu complement
\(C\hookrightarrow\seg{n'}\) (tota part té complement,
capítol~\ref{cap:valors}), de manera que
\(\seg{n}\cong C+\mathbf1_q+\mathbf1_{\ast'}\), on \(\mathbf1_q\) és la
fibra de \(q\) i \(\mathbf1_{\ast'}\) el sumand nou. Sobre aquesta
descomposició, \(\tau\) actua com:
\[
\tau|_C=\id_C,
\qquad
\tau|_{\mathbf1_q}\colon\mathbf1_q\xrightarrow{\ \sim\ }\mathbf1_{\ast'},
\qquad
\tau|_{\mathbf1_{\ast'}}\colon\mathbf1_{\ast'}\xrightarrow{\ \sim\ }\mathbf1_q,
\]
és a dir, la identitat sobre \(C\) i l'intercanvi dels dos singletons.
Aquesta definició per casos (propietat universal de la suma,
notació~\ref{not:suma}) fa immediat que \(\tau\circ\tau=\id_{\seg{n}}\),
ja que intercanviar dues vegades és la identitat sobre cada sumand; per
tant \(\tau\) és una involució, i en particular un isomorfisme. Aleshores
\(\tau\circ f\colon\seg{s\circ m}\hookrightarrow\seg{n}\) és un
monomorfisme amb \(\tau\circ f\circ\ast=\tau\circ\jmath\circ
q=\ast'\), de manera que recau en el Cas~1; aplicant-hi el Cas~1 obtenim
el monomorfisme \(g\colon\seg{m}\hookrightarrow\seg{n'}\).

En tots dos casos hi ha un monomorfisme
\(g\colon\seg{m}\hookrightarrow\seg{n'}\), és a dir \(\mathcal M_{m,n'}\) és
no inicial. Per la hipòtesi d'inducció \(P(m)\), aplicada a \(n'\), obtenim
\(m\le n'\), i doncs \(s\circ m\le s\circ n'=n\) (pas del successor,
proposició~\ref{prop:ordre-total}). Com que \(n\) era arbitrari, això prova
\(P(s\circ m)\) i tanca la inducció.
\end{prova}

\begin{teorema}[Unicitat del cardinal finit]\label{teo:unicitat-cardinal}
Si \(A\) és finit, la part \(\operatorname{Card}(A)\hookrightarrow\N\) té un
\emph{únic} punt: el cardinal \(n\) amb \(\seg{n}\approx A\) és únic. En
conseqüència, la notació \(|A|\) està ben definida, i dos conjunts finits
són isomorfs si i només si tenen el mateix cardinal.
\end{teorema}

\begin{prova}
Treballem sobre els segments interns (observació~\ref{obs:seg-vs-corxet}).
Si \(\seg{m}\approx A\) i \(\seg{n}\approx A\), aleshores
\(\seg{m}\approx\seg{n}\) per transitivitat de l'equipotència
(definició~\ref{def:equipotencia}), és a dir hi ha un isomorfisme
\(\seg{m}\cong\seg{n}\); pel lema~\ref{lema:comparacio-ordinals}, \(m=n\).
Això prova que \(\operatorname{Card}(A)\) té com a molt un punt, i dona la
bona definició de \(|A|\). Per a la darrera afirmació, si \(A\cong B\) amb
\(\seg{|A|}\approx A\) i \(\seg{|B|}\approx B\), aleshores
\(\seg{|A|}\approx\seg{|B|}\) i \(|A|=|B|\); recíprocament, \(|A|=|B|\)
dona \(A\approx\seg{|A|}=\seg{|B|}\approx B\), és a dir \(A\cong B\).
\end{prova}

\section{El principi dels calaixos}\label{sec:calaixos}

La comparació d'ordinals conté ja, essencialment, el principi dels
calaixos. L'enunciem en la forma habitual i en derivem l'equivalència, per
a conjunts finits del mateix cardinal, entre injectivitat, exhaustivitat i
bijectivitat.

\begin{teorema}[Principi dels calaixos]\label{teo:calaixos}
Siguin \(A\) i \(B\) conjunts finits amb \(|A|>|B|\). Aleshores no hi ha
cap monomorfisme \(A\hookrightarrow B\). Dit altrament: tota aplicació
d'un conjunt finit cap a un de cardinal estrictament menor identifica
almenys dos elements.
\end{teorema}

\begin{prova}
Un monomorfisme \(A\hookrightarrow B\) donaria, component amb els
isomorfismes \(\seg{|A|}\approx A\) i \(B\approx\seg{|B|}\)
(definició~\ref{def:conjunt-finit}), un monomorfisme
\(\seg{|A|}\hookrightarrow\seg{|B|}\); pel
lema~\ref{lema:comparacio-ordinals} seria \(|A|\le|B|\), contra la
hipòtesi. La reformulació és la negació de la definició de monomorfisme
(capítol~\ref{cap:aplicacions}): si l'aplicació no és un monomorfisme,
dos elements generalitzats distints tenen la mateixa imatge.
\end{prova}

\begin{observacio}[Les induccions dels recomptes són internes]\label{obs:recomptes-interns}
Les proves que segueixen fan inducció sobre un punt \(n\) en afirmacions
com \enquote{\(A+\seg{n}\) és finit amb cardinal \(|A|+n\)}, és a dir,
\enquote{existeix un isomorfisme \(\seg{|A|+n}\to A+\seg{n}\)}, en què
tant el domini com el codomini varien amb \(n\). Perquè el
teorema~\ref{teo:induccio} s'hi pugui aplicar cal que la propietat sigui
una part de \(\N\); i ho és, per la mateixa tècnica de la
definició~\ref{def:card-A} i de la definició~\ref{def:mono-segments}:
totes les aplicacions \(\seg{j}\to A+\seg{n}\) s'estenen a aplicacions
\(\N\to(A+\N)+\mathbf1\), de manera que \enquote{existeix un isomorfisme}
és la no-inicialitat d'una fibra d'una part d'un exponencial fix, i és
una fórmula interna en \(n\). No repetirem la codificació a cada
recompte; el lector la pot reconstruir seguint el model de
\(\mathcal E_A\).
\end{observacio}

\begin{lema}[Recompte de la suma finita]\label{lema:recompte-suma}
Siguin \(A\) i \(B\) conjunts finits. Aleshores la suma \(A+B\) és finita i
\[
|A+B|=|A|+|B|,
\]
on el segon membre és la suma de naturals (definició~\ref{def:suma-naturals}).
En particular, si \(B\) no és inicial ---és a dir, \(|B|\ge 1\)--- aleshores
\(|A+B|>|A|\).
\end{lema}

\begin{prova}
Com que \(B\approx\seg{|B|}\), n'hi ha prou de provar, per inducció sobre
el punt \(n\) (teorema~\ref{teo:induccio}) i amb \(A\) fixat, que
\(A+\seg{n}\) és finit amb cardinal \(|A|+n\); la propietat és interna en
\(n\) per l'observació~\ref{obs:recomptes-interns}. Escrivim les dues
etapes en termes de \(B\). Per a \(|B|=0\), \(B\approx\seg{0}\cong\mathbf0\)
i la neutralitat del conjunt inicial respecte de la suma
(capítol~\ref{cap:valors}) dona \(A+B\cong A\), d'on
\(|A+B|=|A|=|A|+0\).

Pas de \(|B|=n\) a \(|B|=s\circ n\). Sigui \(B\approx\seg{s\circ
n}\cong\seg{n}+\mathbf1\) (proposició~\ref{prop:segment-successor}). Per
l'associativitat de la suma de conjunts (capítol~\ref{cap:valors}),
\[
A+B\cong A+(\seg{n}+\mathbf1)\cong (A+\seg{n})+\mathbf1 .
\]
Per la hipòtesi d'inducció, \(A+\seg{n}\) és finit amb
\(|A+\seg{n}|=|A|+n\); i afegir un sumand terminal incrementa el cardinal
en un, perquè \((A+\seg{n})+\mathbf1\approx\seg{|A|+n}+\mathbf1\cong
\seg{s\circ(|A|+n)}\) per la descomposició successora
(proposició~\ref{prop:segment-successor}). Així \(A+B\) és finit i
\[
|A+B|=s\circ(|A|+n)=|A|+(s\circ n)=|A|+|B|,
\]
on la penúltima igualtat és la definició recursiva de la suma
(definició~\ref{def:suma-naturals}). Finalment, si \(|B|\ge1\) aleshores
\(|B|=s\circ k\) per a algun \(k\) (lema~\ref{lema:zero-o-successor}) i
\(|A|+|B|=|A|+(s\circ k)\), que és \(>|A|\) per la definició de l'ordre
estricte (definició~\ref{def:ordre-naturals}).
\end{prova}

\begin{lema}[Tota part d'un conjunt finit és finita]\label{lema:part-finita}
Sigui \(S\hookrightarrow A\) una part d'un conjunt finit \(A\). Aleshores
\(S\) i el seu complement són finits, i \(|S|\le|A|\).
\end{lema}

\begin{prova}
Com que \(A\) és finit, \(A\approx\seg{n}\) amb \(n=|A|\); una part de
\(A\) es correspon, per aquest isomorfisme, amb una part de \(\seg{n}\),
així que n'hi ha prou de provar-ho per a \(A=\seg{n}\), per inducció sobre
\(n\) (teorema~\ref{teo:induccio}).

Per a \(n=0\), \(\seg{0}\cong\mathbf0\) i l'única part és la inicial, que
és finita amb cardinal \(0\).

Pas de \(n\) a \(s\circ n\). Per la descomposició successora,
\(\seg{s\circ n}\cong\seg{n}+\mathbf1\)
(proposició~\ref{prop:segment-successor}). Sigui \(S\hookrightarrow
\seg{s\circ n}\) una part. La imatge inversa de \(S\) per les dues
injeccions la descompon en dues parts complementàries: \(S_0\hookrightarrow
\seg{n}\), la part de \(S\) situada sobre \(\seg{n}\), i \(S_1\hookrightarrow
\mathbf1\), la part situada sobre el punt nou. Ara \(S\cong S_0+S_1\)
(les dues parts, complementàries dins \(S\), formen una suma,
capítol~\ref{cap:valors}). Per la hipòtesi d'inducció, \(S_0\) és finit amb
\(|S_0|\le n\). I \(S_1\hookrightarrow\mathbf1\) és inicial o terminal
---les úniques parts de \(\mathbf1\) (capítol~\ref{cap:valors})---, en tot
cas finit amb \(|S_1|\le1\). Pel recompte de la suma
(lema~\ref{lema:recompte-suma}), \(S\) és finit i
\[
|S|=|S_0|+|S_1|\le n+1=s\circ n=|\seg{s\circ n}| .
\]
El complement de \(S\) dins \(\seg{s\circ n}\) es tracta igual, canviant
\(S_0,S_1\) pels seus complements dins \(\seg{n}\) i \(\mathbf1\)
respectivament.
\end{prova}

\begin{teorema}[Injectiu, exhaustiu i bijectiu en el cas finit]\label{teo:endo-finit}
Sigui \(f\colon A\to B\) una aplicació entre conjunts finits amb
\(|A|=|B|\). Són equivalents:
\begin{enumerate}
\item \(f\) és un monomorfisme;
\item \(f\) és un epimorfisme;
\item \(f\) és un isomorfisme.
\end{enumerate}
\end{teorema}

\begin{prova}
Que (3) implica (1) i (2) és immediat. Vegem que (1) implica (3) i,
després, que (2) implica (3).

\emph{(1) \(\Rightarrow\) (3).} Suposem \(f\) un monomorfisme. Per la
factorització imatge (teorema~\ref{teo:factoritzacio}), \(f=m\circ e\) amb
\(e\colon A\twoheadrightarrow\operatorname{Im}(f)\) epimorfisme regular i
\(m\colon\operatorname{Im}(f)\hookrightarrow B\) monomorfisme; com que
\(f\) és un monomorfisme, \(e\) és un isomorfisme
(proposició~\ref{prop:mono-nucli-trivial}), i per tant
\(\operatorname{Im}(f)\cong A\) és finit amb
\(|\operatorname{Im}(f)|=|A|=|B|\). Ara \(m\) exhibeix
\(\operatorname{Im}(f)\) com una part de \(B\). Si \(m\) no fos exhaustiu,
el seu complement \(C\hookrightarrow B\) seria una part no inicial, finita
pel lema~\ref{lema:part-finita} (és una part del finit \(B\)), i \(B\) es
descompondria en la suma disjunta \(B\cong\operatorname{Im}(f)+C\)
(el complement d'una part i la seva part formen una suma,
capítol~\ref{cap:valors}). Com que tots dos sumands són finits, s'hi aplica
el recompte de la suma (lema~\ref{lema:recompte-suma}); i com que \(C\) no
és inicial, \(|B|=|\operatorname{Im}(f)|+|C|>|\operatorname{Im}(f)|=|B|\),
una contradicció. Per tant \(m\) és exhaustiu, i sent un monomorfisme
exhaustiu és un isomorfisme (proposició~\ref{prop:epi-imatge-total} i
corol·lari~\ref{cor:bimorfisme-iso}). Aleshores \(f=m\circ e\) és
composició de dos isomorfismes, i és un isomorfisme.

\emph{(2) \(\Rightarrow\) (3).} Suposem ara \(f\) un epimorfisme. Per
l'axioma d'elecció (axioma~\ref{ax:eleccio}), \(f\) és escindit: existeix
una secció
\(s\colon B\to A\) amb \(f\circ s=\id_B\). Tota secció és un monomorfisme
(proposició~\ref{prop:seccio-mono-retraccio-epi}), de manera que
\(s\colon B\to A\) és un monomorfisme entre conjunts finits amb
\(|B|=|A|\). Per la implicació (1)\(\Rightarrow\)(3) ja provada, \(s\) és un
isomorfisme. Aleshores, de \(f\circ s=\id_B\) se segueix \(f=s^{-1}\)
---component per la dreta amb \(s^{-1}\)---, i \(f\) és un isomorfisme.
Aquesta prova evita tota descomposició en fibres i no és circular: només
usa l'elecció, ja disponible, i la direcció (1)\(\Rightarrow\)(3). És
l'únic lloc del capítol on intervé l'elecció, i s'hi pot evitar: com que
\(A\approx\seg{n}\) és ben ordenat, una secció de \(f\) es pot
\emph{definir} enviant cada punt \(b\) al mínim de la fibra \(f^{-1}(b)\),
que és no buida perquè \(f\) és exhaustiu sobre punts; la construcció
interna d'aquesta secció és un exercici.
\end{prova}

\begin{observacio}[Lectura clàssica]\label{obs:calaixos-classica}
En conjunts, el teorema~\ref{teo:calaixos} és l'enunciat familiar segons
el qual, si es reparteixen \(|A|\) objectes en \(|B|<|A|\) calaixos, algun
calaix en rep més d'un. El teorema~\ref{teo:endo-finit} és el fet, també
familiar i propi \emph{només} del cas finit, que una aplicació d'un
conjunt finit en si mateix és injectiva exactament quan és exhaustiva.
Sobre conjunts infinits ---\(\N\) mateix, amb el successor--- això falla:
el successor és injectiu però no exhaustiu.
\end{observacio}

\section{Parts d'un conjunt finit}\label{sec:parts-finites}

El pas de \(A\) a \(A+\mathbf1\) duplica el nombre de parts: cada part de
\(A+\mathbf1\) o bé conté el punt nou o bé no. Aquesta observació,
convertida en inducció, dona el recompte de parts. Abans, però, ens cal el
recompte del producte, que la prova utilitza sota la forma de la
duplicació \(|A\times\mathbf2|=2\cdot|A|\).

\begin{lema}[Recompte del producte finit]\label{lema:recompte-producte}
Siguin \(A\) i \(B\) conjunts finits. Aleshores el producte \(A\times B\)
és finit i
\[
|A\times B|=|A|\cdot|B|,
\]
on el segon membre és el producte de naturals
(definició~\ref{def:producte-exp}). En particular
\(|A\times\mathbf2|=2\cdot|A|\).
\end{lema}

\begin{prova}
Fem inducció sobre \(|B|\), amb \(A\) fixat (teorema~\ref{teo:induccio}).
Per a \(|B|=0\), \(B\approx\seg{0}\cong\mathbf0\) i
\(A\times\mathbf0\cong\mathbf0\) ---el producte per l'inicial és inicial,
per la distributivitat del producte sobre la suma i l'estrictesa de
\(\mathbf0\) (capítol~\ref{cap:valors})---, d'on \(|A\times B|=0=|A|\cdot0\).

Pas de \(|B|=n\) a \(|B|=s\circ n\). Sigui \(B\approx\seg{s\circ
n}\cong\seg{n}+\mathbf1\) (proposició~\ref{prop:segment-successor}). Per la
distributivitat del producte sobre la suma
(teorema~\ref{teo:distributivitat}),
\[
A\times B\cong A\times(\seg{n}+\mathbf1)\cong (A\times\seg{n})+(A\times\mathbf1)
\cong (A\times\seg{n})+A,
\]
usant \(A\times\mathbf1\cong A\). Per la hipòtesi d'inducció
\(A\times\seg{n}\) és finit amb \(|A\times\seg{n}|=|A|\cdot n\); pel
recompte de la suma (lema~\ref{lema:recompte-suma}),
\[
|A\times B|=|A\times\seg{n}|+|A|=|A|\cdot n+|A|=|A|\cdot(s\circ n)=|A|\cdot|B|,
\]
on la penúltima igualtat és la definició recursiva del producte
(definició~\ref{def:producte-exp}). El cas \(B=\mathbf2\) (amb
\(|\mathbf2|=2\)) dona \(|A\times\mathbf2|=|A|\cdot2=2\cdot|A|\) per la
commutativitat del producte (proposició~\ref{prop:lleis-aritmetica}).
\end{prova}

\begin{teorema}[Nombre de parts d'un conjunt finit]\label{teo:nombre-parts}
Si \(A\) és finit amb \(|A|=n\), aleshores el conjunt de parts
\(\mathcal P(A)\) (definició~\ref{def:conjunt-parts}) és finit i
\[
|\mathcal P(A)|=2^{n},
\]
on \(2^{n}\) és l'exponenciació de naturals
(definició~\ref{def:producte-exp}).
\end{teorema}

\begin{prova}
Com que \(\mathcal P(A)\cong\mathbf2^{A}\)
(definició~\ref{def:conjunt-parts}) i \(A\approx\seg{n}\) amb \(n=|A|\),
n'hi ha prou de provar \(|\mathbf2^{\seg{n}}|=2^{n}\) per inducció sobre
\(n\). Per a \(n=0\), \(\mathbf2^{\seg{0}}\cong\mathbf2^{\mathbf0}\cong
\mathbf1\), i \(2^{0}=1\) (definició~\ref{def:producte-exp}). Pas de \(n\)
a \(s\circ n\): de \(\seg{s\circ n}\cong\seg{n}+\mathbf1\)
(proposició~\ref{prop:segment-successor}) i de la conversió d'un
exponencial d'una suma en un producte d'exponencials
(proposició~\ref{prop:exponencial-suma}),
\[
\mathbf2^{\seg{s\circ n}}\cong\mathbf2^{\seg{n}+\mathbf1}\cong
\mathbf2^{\seg{n}}\times\mathbf2^{\mathbf1}\cong\mathbf2^{\seg{n}}\times\mathbf2 .
\]
Per la hipòtesi d'inducció \(\mathbf2^{\seg{n}}\) és finit de cardinal
\(2^{n}\), i multiplicar per \(\mathbf2\) duplica el cardinal
(lema~\ref{lema:recompte-producte}), de manera que
\(\mathbf2^{\seg{s\circ n}}\) té cardinal \(2^{n}+2^{n}=2\cdot2^{n}=
2^{s\circ n}\), on la darrera igualtat és la definició recursiva de
l'exponenciació.
\end{prova}

\begin{observacio}[Els recomptes de la suma i del producte]\label{obs:recompte-producte-suma}
Els dos recomptes que aquestes proves fan servir ---\(|A+B|=|A|+|B|\) i
\(|A\times B|=|A|\cdot|B|\) per a conjunts finits--- s'han establert com a
lemes previs (lemes~\ref{lema:recompte-suma} i~\ref{lema:recompte-producte}),
tots dos per inducció sobre el cardinal del segon factor. Són la traducció,
al recompte de cardinals, de la suma disjunta i de l'enumeració fila per
fila del producte; la duplicació \(|A\times\mathbf2|=2\cdot|A|\) que hem
usat n'és el cas \(B=\mathbf2\).
\end{observacio}

\section{Recompte d'aplicacions}\label{sec:recompte-aplicacions}

L'exponencial \(B^{A}\) internalitza el conjunt de les aplicacions de
\(A\) cap a \(B\) (capítol~\ref{cap:exponencials}). Quan tots dos conjunts
són finits, en compta exactament les aplicacions.

\begin{teorema}[Nombre d'aplicacions entre conjunts finits]\label{teo:nombre-aplicacions}
Si \(A\) i \(B\) són finits amb \(|A|=n\) i \(|B|=m\), aleshores
l'exponencial \(B^{A}\) és finit i
\[
|B^{A}|=m^{n} .
\]
\end{teorema}

\begin{prova}
Fixats isomorfismes \(A\approx\seg{n}\) i \(B\approx\seg{m}\) (amb
\(n=|A|\), \(m=|B|\)), l'exponencial només depèn dels dos conjunts llevat
d'isomorfisme, de manera que n'hi ha prou de provar
\(|\seg{m}^{\seg{n}}|=m^{n}\) per inducció sobre \(n\), uniformement en
\(m\). Per a \(n=0\), \(\seg{m}^{\seg{0}}\cong\seg{m}^{\mathbf0}\cong
\mathbf1\), i \(m^{0}=1\). Pas de \(n\) a \(s\circ n\): de
\(\seg{s\circ n}\cong\seg{n}+\mathbf1\)
(proposició~\ref{prop:segment-successor}) i de la conversió de
l'exponencial d'una suma (proposició~\ref{prop:exponencial-suma}),
\[
\seg{m}^{\seg{s\circ n}}\cong\seg{m}^{\seg{n}+\mathbf1}\cong
\seg{m}^{\seg{n}}\times\seg{m}^{\mathbf1}\cong\seg{m}^{\seg{n}}\times\seg{m} .
\]
Per la hipòtesi d'inducció \(\seg{m}^{\seg{n}}\) té cardinal \(m^{n}\), i
pel recompte del producte (lema~\ref{lema:recompte-producte}) multiplicar
per \(\seg{m}\) multiplica el cardinal per \(m\); així
\(|\seg{m}^{\seg{s\circ n}}|=m^{n}\cdot m=m^{s\circ n}\), la definició
recursiva de l'exponenciació.
\end{prova}

\begin{observacio}[Lectura clàssica i enllaç amb la pràctica]\label{obs:recompte-classica}
Els quatre recomptes ---\(2^{n}\) parts, \(m^{n}\) aplicacions,
\(|A+B|=|A|+|B|\) i \(|A\times B|=|A|\cdot|B|\)--- són l'aritmètica
combinatòria elemental, ara obtinguda sense sortir de la teoria: cada
igualtat és un isomorfisme construït per inducció, no un recompte de
punts. Aquests enunciats són el punt de partida dels exercicis
computacionals de la part pràctica, on es calculen cardinals concrets i
es refà, en casos petits, la inducció general. En particular, el nombre
d'aplicacions injectives i el de bijeccions de \([k]\) en si mateix
---les permutacions--- es tracten allà com a aplicació directa del
principi dels calaixos i del teorema~\ref{teo:endo-finit}.
\end{observacio}

\begin{resumcap}
\apartat{Idees centrals}
\begin{enumerate}
  \item Dues famílies de \enquote{nombres finits}: els ordinals \([k]=\mathbf1+\dots+\mathbf1\), definits al metanivell per a numerals estàndard, i els segments \(\seg{n}=\{k<n\}\), fibres de \(\operatorname{Lt}\hookrightarrow\N\times\N\), que formen una família interna. Coincideixen externament sobre numerals; la finitud es defineix amb els segments.
  \item \(\operatorname{Card}(A)\hookrightarrow\N\) és la imatge de la projecció d'una part \(\mathcal E_A\) d'un exponencial fix que codifica les bijeccions \(\seg{n}\to A\); \(A\) és finit si \(\operatorname{Card}(A)\) té algun punt, i el cardinal és aquest punt, únic.
  \item La unicitat ve de la comparació de segments: un monomorfisme \(\seg{m}\hookrightarrow\seg{n}\) força \(m\le n\). La prova és una inducció interna sobre un predicat codificat, amb el truc de permutar dos punts per reduir al cas en què el punt nou va al punt nou.
  \item Principi dels calaixos; en cardinal igual, mono, epi i iso coincideixen. L'elecció només s'usa per escindir l'epimorfisme, i es pot evitar amb el bon ordre.
  \item Recomptes per inducció sobre el cardinal: \(|A+B|=|A|+|B|\), \(|A\times B|=|A|\cdot|B|\), \(|P(A)|=2^{|A|}\), \(|B^A|=|B|^{|A|}\); cadascun és un isomorfisme construït amb la descomposició successora, la distributivitat i l'exponencial d'una suma.
  \item Les induccions sobre afirmacions d'existència d'isomorfismes són internes perquè les aplicacions entre objectes variables es codifiquen dins exponencials fixos.
\end{enumerate}
\apartat{Errors típics} Definir la finitud com \enquote{isomorf a algun \([k]\)}, que és un esquema extern i no una propietat interna; fer inducció sobre \enquote{existeix un isomorfisme \(\seg{n}\cong\cdots\)} sense comprovar que és una part de \(\N\); creure que \(\operatorname{Card}(A)\) necessita triar un isomorfisme (es construeix com una imatge, sense cap tria); oblidar que \enquote{un punt d'una suma ve d'un sumand} val sobre punts globals i no sobre elements generalitzats, cosa essencial al lema de comparació; i pensar que mono \(\Rightarrow\) iso en cardinal igual val per a conjunts infinits (el successor de \(\N\) ho desmenteix).
\apartat{Lectura recomanada} \cite[§§7.6--7.7 i cap.~8]{lawvere-rosebrugh} per als conjunts finits, el principi dels calaixos i els recomptes en conjunts abstractes; \cite[article~V, sessió~30]{lawvere-schanuel} per al recompte d'aplicacions \(|B^A|=|B|^{|A|}\) amb exemples; \cite[§§12--13]{halmos} per als nombres naturals finits i l'aritmètica de cardinals finits en la teoria clàssica; i \cite[§§7.1--7.2]{awodey} per a l'objecte de nombres naturals i la seva propietat universal en una categoria, que és el que fa possible la codificació interna d'aquest capítol.
\end{resumcap}
